

\section{Introduction}

%"The Moravec Paradox" is currently one of the significant challenges remaining in artificial intelligence technology \cite{Moravec1988mind}. This paradox refers to the contradiction where tasks that even children can perform semi-unconsciously are exceedingly difficult for the latest artificial intelligence and robots compared to complex logical reasoning, such as puzzles or chess. Tasks involving independent movement like walking or running made significant technological advancements in the 2000s \cite{Garcia1996,Hirai1998}, while perception tasks such as vision and voice recognition saw substantial progress in the 2010s through deep learning techniques~\cite{Krizhevsky2012nips,He2016cvpr,Ren2016pami,Redmon2016cvpr,Hinton2013spm,Graves2014icml}. However, performing a variety of tasks using a single common hand, just like humans, remains extremely challenging even with the latest technology available today.
"The Moravec Paradox" is currently one of the significant challenges in artificial intelligence technology \cite{Moravec1988mind}. This paradox refers to the contradiction in which tasks that even children can perform semi-unconsciously are exceedingly difficult for the latest artificial intelligence and robots compared to complex logical reasoning tasks, such as puzzles or chess. Tasks involving independent movements such as walking or running made significant technological advancements in the 2000s \cite{Garcia1996,Hirai1998}. In contrast, perception tasks such as vision and voice recognition saw substantial progress in the 2010s through deep learning techniques~\cite{Krizhevsky2012nips,He2016cvpr,Ren2016pami,Redmon2016cvpr,Hinton2013spm,Graves2014icml}. However, performing a variety of tasks using a single common hand, just like humans, remains extremely challenging even with the latest technology available today.



\subsection{Experience-based Robotics}

%The use of deep learning methods is expanding in the fields of both environmental recognition and motion generation.  Examples of this include extracting graspable regions of objects~\cite{Pinto2016icra,Mahler2017dex,Mousavian2019iccv} and combinatorial exploration of action patterns through reinforcement learning~\cite{Levine2016jmlr,Levine2017ijrr}. These research examples represent an inductive approach to directly acquiring models of the environment and trajectory generation from data. In the context of robot learning, data can be generated or acquired as a result of its bottom-up exploratory actions. Sensory-motor information obtained from active interaction, or "experience," plays an important role in robot learning. For example, while conventional learning data is annotated with labels provided by humans, the experience can be utilized in the robot's own motor skill learning even in the absence of teacher labels. Another important perspective that emerges in this concept called experience-based robotics is the functional separation of environmental perception and motion generation. Considering the now common end-to-end approach in processing images, audio, and natural language using deep learning, it becomes evident that the clear separation of these two functions could also be a subject of learning.
Deep learning methods are expanding in the environmental recognition and motion generation fields. Examples include the extraction of graspable regions of objects~\cite{Pinto2016icra,Mahler2017dex,Mousavian2019iccv} and the combinatorial exploration of action patterns through reinforcement learning~\cite{Levine2016jmlr,Levine2017ijrr}. These research examples represent an inductive approach for directly acquiring models of the environment and trajectory generation from data. In robot learning, data can be generated or acquired through bottom-up exploratory actions. Sensory-motor information obtained from active interaction, or "experience," plays an important role in robot learning. For example, while conventional learning data are annotated with labels provided by humans, experience can be utilized in a robot's motor-skill learning, even in the absence of teacher labels. Another important perspective from this concept, called experience-based robotics, is the functional separation between environmental perception and motion generation. Considering the common end-to-end approach for processing images, audio, and natural language using deep learning, it is evident that a clear separation of these two functions could also be a subject of learning.


%The concept of recognizing and generating partial "functions" of autonomous systems here is important, as it is believed to emerge gradually from the interaction of the body, environment, neural circuits, and the continuous dynamics of end-to-end learning, as proposed in cognitive developmental robotics~\cite{Asada2015} or brain-inspired learning~\cite{Haruno2001}. Particularly, to deploy robots in the real world, it is crucial to bridge the gap between the motion model and reality with limited data. While deep learning-based approaches show improvements in task accuracy through large amounts of data and trials, data collection for model training is highly costly. In particular, robot motion data depends on hardware, making it uncommon to have access to extensive data for model tuning. Although research is conducted to transfer learning results from simulation environments to the real world, modeling nonlinear physical phenomena remains challenging and is not entirely solvable. It is necessary to achieve a wide range of task behaviors by creatively combining representations derived from the limited motion experiences of robots.
The concept of recognizing and generating partial "functions" of autonomous systems here is important, as it is believed to emerge gradually from the interaction of the body, environment, neural circuits, and the continuous dynamics of end-to-end learning, as proposed in cognitive developmental robotics~\cite{Asada2015} or brain-inspired learning~\cite{Haruno2001}. In particular, to deploy robots in the real world, it is crucial to bridge the gap between motion models and reality using limited data. Although deep learning-based approaches show improvements in task accuracy through large amounts of data and trials, data collection for model training is costly. In particular, robot motion data depend on the hardware, making accessing extensive data for model tuning uncommon. Although research has been conducted to transfer learning results from simulation environments to the real world, modeling nonlinear physical phenomena remains challenging and unsolvable. It is necessary to achieve a wide range of task behaviors by creatively combining representations derived from the limited motion experience of robots.



\subsection{Proposed Robot Learning Concept}

%In this paper, we discuss the concept of "deep predictive learning," a methodology for applying deep learning techniques to robotics, along with its implementation methods. This concept, inspired by predictive coding theory~\cite{Rao1999, Friston2006jpp, Friston2010nrn}, is based on the fundamental strategy of predicting the near-future sensory-motor states of robots and online minimization of the prediction error between the real world and the model. Furthermore, by embedding the temporal relationships between sensory-motor information in systems where the environment and body interact dynamically, the system transitions to other motion dynamics based on prediction error during generation, enabling the execution of high-dimensional, long-term tasks.
This study discusses the concept of "deep predictive learning," a methodology for applying deep learning techniques to robotics, along with its implementation methods. This concept, inspired by predictive coding theory~\cite{Rao1999, Friston2006jpp, Friston2010nrn}, is based on the fundamental strategy of predicting the near-future sensory-motor states of robots and online minimization of the prediction error between the real world and the model. Furthermore, by embedding the temporal relationships between sensory and motor information in systems where the environment and body interact dynamically, the system transitions to other motion dynamics based on prediction errors during generation, thereby enabling the execution of high-dimensional, long-term tasks.


%In Section 2 of this paper, we discuss previous research and its challenges related to the application of deep learning in robotics and describe the positioning of deep predictive learning. Section 3 covers the background of deep predictive learning, including predictive coding and the free energy principle, and outlines strategies for applying them to robot motion learning. We also discuss specific implementation approaches based on deep learning techniques and examine them from the perspective of cognitive robotics, including dynamic cognitive systems and action-normative systems. In Section 4, we summarize the published studies that we have developed using the DPL concept from two perspectives: scalability by combining multiple motions, and application to complex manipulation tasks. Finally, in Section 5, we provide a summary and discuss future prospects.
Section 2 discusses previous research and the challenges related to applying deep learning in robotics and describes the positioning of deep predictive learning. Section 3 covers the background of deep predictive learning, including predictive coding and the free-energy principle, and outlines strategies for applying them to robot motion learning. We also discuss specific implementation approaches based on deep learning techniques and examine them from the perspective of cognitive robotics, including dynamic cognitive and action-normative systems. Section 4 summarizes the published studies we developed using the proposed concept from two perspectives: scalability by combining multiple motions and application to complex manipulation tasks. Finally, in Section 5, we summarize and discuss prospects.

